\hl{This chapter fixes what the system must do, the resources it needs,
and the societal, environmental, ethical, economic and practical
constraints under which it was designed and evaluated. The framework
answers questions written in Standard Bangla, Chittagonian or Sylheti,
in Bangla script or romanized form, about three domains (school
education, general knowledge and healthcare), and always responds in
Standard Bangla.}

\section{Final Specifications and Requirements}\label{final-specifications-and-requirements}

\hl{The final system routes each question by language variety and
domain, normalizes dialectal and romanized questions into Standard
Bangla, and answers them with domain-specific retrieval-augmented
generation. It is designed to give grounded Standard Bangla answers to
users of Standard Bangla, Chittagonian and Sylheti, and to ask for
clarification, abstain or refer the user to a doctor when an answer
would not be reliable.}

\subsection{Functional Requirements}\label{functional-requirements}

\hl{The final system satisfies the following functional requirements
(FR).}

\begin{longtable}{@{}p{0.080\textwidth}p{0.370\textwidth}p{0.220\textwidth}p{0.220\textwidth}@{}}
\caption{Functional requirements, implementing components and verifying experiments.}\label{tab:functional-requirements}\\
\toprule
\textbf{ID} & \textbf{Requirement} & \textbf{Implemented by} & \textbf{Verified by} \\
\midrule
\endfirsthead
\toprule
\textbf{ID} & \textbf{Requirement} & \textbf{Implemented by} & \textbf{Verified by} \\
\midrule
\endhead
\bottomrule
\endfoot
FR1 & Accept questions in Standard Bangla, Chittagonian and Sylheti, in Bangla script or romanized & script detection + dialect router & router test set, end-to-end test \\
FR2 & Identify the language variety of each question & calibrated TF-IDF dialect router & per-class F1, unseen dialects \\
FR3 & Rewrite dialectal or romanized questions into Standard Bangla, preserving symptoms, numbers and named entities & LoRA normalization expert & chrF++/BLEU on held-out pairs \\
FR4 & Decide the domain (education, general, medical, out-of-scope) & domain router & per-class F1, end-to-end domain accuracy \\
FR5 & Ground every answer in retrieved evidence from the domain knowledge base & per-domain retriever (BM25 or dense) & hit@k, MRR, nDCG@10 \\
FR6 & Answer with a domain-specific expert on one shared base model & Mixture-of-Adapters (LoRA experts) & F1 / chrF++ per domain, pooled-adapter ablation \\
FR7 & Answer, ask for clarification, or abstain when information is missing or the question is out of scope & clarification/\allowbreak{}abstention controller & controller test set, end-to-end behaviour table \\
FR8 & Never give emergency, dosage or prescription advice; always refer to a registered doctor & rule-based medical safety layer & safety test set, referral rate \\
FR9 & Respond in Standard Bangla with a transparent trace (variety, normalized question, domain, decision, evidence) & pipeline output & chatbot demonstration \\
FR10 & Produce every table, figure and statistical test of the thesis reproducibly, reported separately for each variety & staged notebooks + manifests & regenerated result files \\
\bottomrule
\end{longtable}

\subsection{Non-Functional Requirements}\label{non-functional-requirements}

\begin{itemize}
\item \hl{\textbf{Fairness across varieties:} Every result is reported
separately for Standard Bangla, Chittagonian and Sylheti, and the
Standard-vs-dialect gap is a primary metric.}

\item \hl{\textbf{No harm to Standard input:} Dialect handling must not
reduce quality for Standard Bangla questions (tested by the
NO\_NORMALIZATION and ALWAYS\_NORMALIZE ablations).}

\item \hl{\textbf{Efficiency:} One 4-bit base model serves every domain;
each additional domain costs one adapter of 133 MB instead of another
large language model. End-to-end latency is about 0.57 s per question in
batched evaluation on one GPU.}

\item \hl{\textbf{Modularity:} The routers, normalization expert,
retrievers, knowledge bases and domain adapters can each be replaced or
retrained without retraining the others, so that new dialects or domains
can be added later; one domain expert trains in 55--113 minutes.}

\item \hl{\textbf{Reproducibility:} Fixed seeds (17; 42 and 101 for
finalists), pinned software versions, configuration files, SHA-256
manifests of every data transfer, and model selection on validation data
only.}

\item \hl{\textbf{Robustness of execution:} Every GPU job is resumable from
checkpoints, so interruption of a cloud machine never loses finished
work.}

\item \hl{\textbf{Privacy:} No user data are stored; the knowledge bases are
public resources and patient questions never leave the system (no
external API).}
\end{itemize}

\subsection{Technical and Resource Requirements}\label{technical-and-resource-requirements}

\begin{itemize}
\item \hl{\textbf{Software:} Python 3.12, PyTorch 2.8.0 (CUDA 12.8), Hugging
Face transformers 4.57.6 (and 5.17.0 in a separate environment for
Gemma-4), PEFT 0.21.0, bitsandbytes 0.50.2, accelerate 1.15.0,
scikit-learn 1.6.1 and sacreBLEU 2.5.1.}

\item \hl{\textbf{Retrieval stack:} BM25 for sparse retrieval and BGE-M3 and
multilingual-E5 embeddings for dense retrieval \cite{robertson2009bm25}, \cite{chen2024bgem3},
\cite{wang2024me5}, with BanglaBERT as the transformer baseline for the routers
\cite{bhattacharjee2022banglabert}.}

\item \hl{\textbf{CPU compute:} Google Colab CPU runtimes (browser or VS
Code) for data preparation, routers, sparse retrieval, analysis and
report generation.}

\item \hl{\textbf{GPU compute:} One NVIDIA A40 (48 GB, 44.4 GB usable)
rented on RunPod for model comparison, expert training, dense retrieval
and full-system evaluation. All language-model training uses 4-bit
quantization and LoRA adapters so that it fits in a single GPU
\cite{dettmers2023qlora}.}

\item \hl{\textbf{Storage and transfer:} Google Drive for data and results,
a private Hugging Face dataset repository for moving stage outputs
between Colab and RunPod, and an 80 GB volume on the GPU machine.}

\item \hl{\textbf{Human resources:} Native speakers of Chittagonian and
Sylheti to verify dialect questions and a physician to review medical
answers. Review sheets for both were prepared, but these reviews could
not be completed within the project time (Section 5.11).}
\end{itemize}

\subsection{Final Design Specification}\label{final-design-specification}

\hl{Table~\ref{tab:final-design} summarises the configuration that was finally selected.}

\begin{longtable}{@{}p{0.270\textwidth}p{0.630\textwidth}@{}}
\caption{Final specification of the selected design.}\label{tab:final-design}\\
\toprule
\textbf{Component} & \textbf{Selected configuration} \\
\midrule
\endfirsthead
\toprule
\textbf{Component} & \textbf{Selected configuration} \\
\midrule
\endhead
\bottomrule
\endfoot
Base language model & Gemma-4-E4B-it \cite{gemma4}, loaded once in 4-bit NF4 with double quantization \\
Experts (Mixture-of-Adapters) & LoRA adapters \cite{hu2022lora} (r = 16, \ensuremath{\alpha} = 32, dropout 0.05, all attention and MLP projections; 34.9 M trainable parameters, 133 MB each): normalization (2 epochs, selected on validation data), education, general, medical; hard top-1 routing \\
Dialect router & char (1--5) + word (1--2) TF-IDF, calibrated linear SVM; test macro-F1 0.946 \\
Domain router & same features and classifier over education / general / medical / other; test macro-F1 0.908 \\
Controller & gradient-boosted trees on router probabilities, retrieval scores and surface features; decisions ANSWER / CLARIFY / ABSTAIN \\
Retrieval & education and general: BM25; medical: dense BGE-M3 \cite{chen2024bgem3}; 3 passages per question \\
Knowledge bases & education 163,591 chunks, general 7,948, medical 28,134 \\
Safety layer & rules for emergencies (fixed reply pointing to the national emergency number 999), dose/prescription refusal, referral note on every medical answer \\
Output & Standard Bangla answer + trace (variety, normalized question, domain, decision, evidence) \\
\bottomrule
\end{longtable}

\section{Societal Impact}\label{societal-impact}

\hl{Chittagonian and Sylheti are spoken by millions of people in
Bangladesh and the diaspora, yet question answering systems and language
models are developed and evaluated predominantly on Standard Bangla. Our
results quantify the consequence: without normalization the same
questions score 33\% (Chittagonian) and 43\% (Sylheti) lower than in
Standard Bangla (Section 5.6). The expected positive impacts are:}

\begin{itemize}
\item \hl{\textbf{Access to health information:} People can describe
symptoms in their own variety and receive evidence-grounded general
health information, with emergencies redirected to 999 and a doctor
referral on every medical answer. This matters in a country where only
about 11\% of doctors and nurses work in rural areas, home to about two
thirds of the population \cite{nuruzzaman2022informing}.}

\item \hl{\textbf{Education support:} Students can ask textbook questions
(national curriculum) in their home variety and receive answers grounded
in the same textbooks they study \cite{nctbqa}.}

\item \hl{\textbf{Digital inclusion:} Speakers who write in romanized Bangla
or in a regional variety are served by the same system, instead of being
expected to switch to the standard written form.}

\item \hl{\textbf{Research infrastructure:} The evaluation protocol (paired
questions in three varieties, per-variety reporting, human-written
dialect test questions) makes the dialect gap measurable and can be
reused for other Bangla varieties, together with the trained routers,
normalization adapter and evaluation scripts.}
\end{itemize}

\hl{The main societal risks are over-reliance on automatic health
answers and wrong answers in dialects the system does not know well.
They are mitigated by the informational-only design, abstention when
evidence is missing, the referral note, and transparent reporting of
failure cases. The system is a research prototype and is not intended
for clinical use without further validation by physicians, regulatory
review and ongoing monitoring.}

\section{Ethical Issues}\label{ethical-issues}

\begin{itemize}
\item \hl{\textbf{Medical safety:} The medical branch is informational only
and the system does not present itself as a doctor. Emergency
descriptions bypass the language model and return a fixed message to
call 999; dose or prescription requests are refused; every medical
answer ends with a referral to a registered doctor. Advice that could
delay urgent care is treated as the most serious type of error. No
language model is used to judge safety; a physician review sheet was
prepared for human validation.}

\item \hl{\textbf{Data licensing and attribution:} Only public research
datasets were used and each is cited. Several (e.g. BanglaRQA \cite{ekram2022banglarqa},
released under CC BY-NC-SA 4.0, and BanglaCHQ-Summ \cite{khan2023banglachqsumm}) carry
non-commercial licences, so the trained system is for research use. The
dialect health dataset BanglaCHQ-Prantik \cite{prantik2025} was not used without
the authors\textquotesingle{} permission.}

\item \hl{\textbf{Privacy:} No personal data were collected. Knowledge-base
texts come from published resources; the demonstration stores no
conversations. Closed commercial models were deliberately not used, so
that health questions never leave the system.}

\item \hl{\textbf{Fairness and representativeness:} Performance is reported
per variety, and the remaining dialect gap is stated explicitly rather
than averaged away. Dialect data are imbalanced across varieties, and
some sources, such as NCTB-QA, contain model-generated questions
\cite{nctbqa}; these limitations are stated with the results.}

\item \hl{\textbf{Respect for dialects:} Normalizing a dialect into Standard
Bangla is a technical step to reuse existing resources. It does not
imply that regional varieties are incorrect or inferior; answering in
the user\textquotesingle s own dialect is listed as future work.}

\item \hl{\textbf{Transparency and trust:} Every answer is returned with a
trace (detected variety, normalized question, domain, decision and
evidence), and the system abstains or asks for clarification when it is
not sure.}

\item \hl{\textbf{Honest reporting:} Negative and neutral findings are
reported: the dialect-invariant encoder did not beat its base encoder,
routed experts do not beat a single pooled adapter on accuracy, and the
Chittagonian result remains below a published specialised model.}
\end{itemize}

\section{Standards}\label{standards}

\hl{The work follows these standards and guidelines (alignment with
their principles; no formal certification is claimed):}

\begin{itemize}
\item \hl{\textbf{Text encoding:} All text is stored in UTF-8 and normalized
to Unicode Normalization Form C (Unicode Standard Annex \#15) with
invisible characters removed, so that identical Bangla strings compare
equal.}

\item \hl{\textbf{Evaluation reporting:} Corpus BLEU and chrF++ are computed
with sacreBLEU \cite{post2018sacrebleu}; confidence intervals and paired significance
tests follow common NLP practice \cite{koehn2004bootstrap}, \cite{riezler2005randomization}; multiple
comparisons are corrected with the Holm procedure \cite{holm1979}.}

\item \hl{\textbf{Reproducibility and data management:} Fixed
train/validation/test splits (official splits preserved), fixed random
seeds, pinned package versions, versioned configuration files, a leakage
audit, and SHA-256 checksums for every transferred file.}

\item \hl{\textbf{Responsible AI for health:} The safety design follows
widely accepted principles for AI in health information, including the
World Health Organization\textquotesingle s guidance \cite{who2021ethics}: the
system gives general information only, defers diagnosis and prescribing
to qualified professionals, and escalates emergencies.}
\end{itemize}

\section{Project Management Plan}\label{project-management-plan}

\hl{The project is organized into phases aligned with the thesis
timeline. The plan prioritizes early data readiness and a working
baseline pipeline, followed by component improvement and evaluation.}

\begin{itemize}
\item \hl{\textbf{Phase 1: Literature review and planning}}

\hl{Finalize the scope, target dialects, domains and evaluation metrics,
and confirm the availability and licenses of the datasets.}

\item \hl{\textbf{Phase 2: Data collection and preprocessing}}

\hl{Collect the dialect, QA and healthcare datasets, clean and normalize
them, build the unified data schema, and construct the three knowledge
bases (education, general knowledge and medical).}

\item \hl{\textbf{Phase 3: Router and normalization expert}}

\hl{Train and evaluate the variety and domain router, and train the
dialect normalization adapter and compare it with zero-shot and
few-shot prompting baselines.}

\item \hl{\textbf{Phase 4: Domain RAG experts and base-model comparison}}

\hl{Build and select the retriever of each domain (BM25 or dense), train
the healthcare, education and general-knowledge adapters, and compare the
candidate base models.}

\item \hl{\textbf{Phase 5: Dialectal test set, evaluation and ablations}}

\hl{Build the paired dialect test set, run the end-to-end and ablation
experiments, and evaluate the system on dialect questions written by
native speakers.}

\item \hl{\textbf{Phase 6: Writing and finalization}}

\hl{Consolidate the results, prepare the final figures and tables, and
complete thesis formatting and proofreading.}
\end{itemize}

\hl{\textbf{Current status.} All six phases have been completed. The data
stage cleaned and split more than twenty public sources into
dialect--Standard pairs, QA items and three knowledge bases (education,
general knowledge and medical), and an automatic leakage audit found no
overlap between training, knowledge-base and test data. The variety and
domain routers and the dialect normalization expert were trained and
evaluated (dialect-router macro-F1 0.946; normalization chrF++ 76.3 for
Chittagonian and 74.0 for Sylheti). Seven open base models were compared
under one QLoRA budget, Gemma-4-E4B-it was selected, and the domain
experts and retrievers were trained and chosen on validation data. The
complete system was evaluated under nine conditions on 2,225 paired test
questions and on 481 dialect questions written by native speakers, and
the thesis is in its final writing stage.}

\section{Risk Management}\label{risk-management}

\hl{The project faced both technical and practical risks. The
mitigations were planned in advance, and Table~\ref{tab:risks} also records what
actually happened.}

\begin{longtable}{@{}L{0.200\textwidth}L{0.360\textwidth}L{0.340\textwidth}@{}}
\caption{Risks, mitigations and what actually happened.}\label{tab:risks}\\
\toprule
\textbf{Risk} & \textbf{Mitigation} & \textbf{Outcome} \\
\midrule
\endfirsthead
\toprule
\textbf{Risk} & \textbf{Mitigation} & \textbf{Outcome} \\
\midrule
\endhead
\bottomrule
\endfoot
Cloud GPU stopped or lost & resumable jobs with atomic checkpoints; results pushed to a private repository after every stage & pods were stopped and replaced several times without losing work \\
Newest model not loadable by the stable software stack & separate environment with transformers 5 for Gemma-4 & Gemma-4 trained and evaluated normally \\
Out-of-memory or silently truncated training data & 4096-token limit after a tokenizer-fertility audit; per-model count of dropped examples & 0 examples dropped in every expert \\
Test data leaking into training or knowledge bases & group-wise splits and an automatic leakage audit & 0 leaks in all six audits \\
Moving to a new cloud account & migration notebook comparing rows read per source before and after & identical data on the new account \\
Native-speaker dialect test data unavailable & human-written question sentences of held-out corpus test splits & 481 human-written test questions evaluated \\
Harmful medical output & rule-based safety layer, referral note, abstention & emergencies handled; remaining gaps reported (Section 5.9) \\
Budget overrun & cost-aware schedule, CPU for all non-GPU work, deleting storage after use & total GPU cost about US\$20 \\
Too little Chittagonian data, especially in the medical domain & combine all available dialect sources; one normalization adapter for both dialects with the variety named in the prompt & 16,030 Chittagonian training pairs; medical Chittagonian remains the scarcest \\
Normalization changes the meaning of a question & confirmation rule for uncertain normalizations; comparison with base-model and prompting baselines & chrF++ 76.3 (Chittagonian) and 74.0 (Sylheti); the confirmation rule never fired, and some dialect words were kept literally \\
Router confuses varieties that share features & calibrated router; every dialect label triggers normalization & macro-F1 0.946; 97.4\% / 94.3\% of human-written Chittagonian / Sylheti recognised, but only 71.5\% of synthetic Chittagonian \\
Routed design does not outperform a single pipeline & pooled adapter with the same number of training steps as a control & equal accuracy (\ensuremath{\pm}0.005); the value of routing is modularity \\
\bottomrule
\end{longtable}

\section{Economic Analysis}\label{economic-analysis}

\hl{Development used free CPU notebooks and about 40 rented A40
GPU-hours at US\$0.49 per hour, i.e. about US\$20 of GPU time plus a few
dollars of storage, well within a student budget. Serving is also
inexpensive: in batched evaluation the complete pipeline answered 2,225
questions in 1,268 s on one A40 (about 0.57 s per question), i.e.
roughly 6,300 questions per GPU-hour, or well under US\$0.001 per
question at the same price. Because all domains share one quantized base
model, adding a domain requires training one adapter (55--113 minutes in
our runs) rather than hosting another model, and an existing domain can
be updated without retraining the others (55 minutes for the medical
expert versus 317 minutes for a single pooled adapter covering all
tasks). Open models also avoid per-token API fees and keep health
questions local. Further economic benefits are:}

\begin{itemize}
\item \hl{\textbf{Reuse instead of new data collection:} Normalizing
dialectal questions lets the system reuse existing Standard Bangla
textbooks, QA datasets and health content instead of building separate
resources for every dialect.}

\item \hl{\textbf{Potential savings for users:} Reliable first-line health
information and curriculum-aligned learning support in a
user\textquotesingle s own variety could reduce unnecessary travel,
delayed care and reliance on paid private tutoring.}
\end{itemize}
